\chapter{Conclusioni e sviluppi futuri}

	In questo lavoro di tesi si è compiuto un confronto tra due due applicazioni realizzate seguendo approcci alla progettazione differenti, da un lato la modalità tradizionale di sviluppo, dall'altra una corrente emergente in crescita di anno in anno.
	
	L'utilizzo di Kubernetes rende più efficiente lo sviluppo e facilita la gestione del sistema, da sottolineare sono in particolari le caratteristiche di \emph{self-healing} e scalabilità offerte.
	
	Gli sviluppi futuri potrebbero essere molteplici e spaziare dall'ambito della sicurezza alla migrazione verso nuove tecnologie.
	L'aspetto sicurezza è indubbiamente di primaria importanza, le applicazioni sviluppate espongono infatti alla possibilità che un utente malevolo riesca ad impossessarsi di un token destinato a un altro utente, utilizzandolo per nascondere la propria identità.
	In secondo luogo la migrazione da PostgreSQL ad un database NoSQL porterebbe a un incremento della velocità di esecuzione delle query e permetterebbe la scalabilità orizzontale del sistema, aspetto di primaria importanza per un sistema distribuito.
	Valide alternative potrebbero essere MongoDB o Cassandra.
	Infine si potrebbe riscrivere l'intera applicazione in un linguaggio compilato: sebbene Python sia un linguaggio moderno, flessibile e sintetico, esso risulta essere uno dei linguaggi più lenti in esecuzione a causa della presenza dell'interprete;
	inoltre si presta bene solo per la scrittura di script o piccole applicazioni, altri linguaggi offrono un supporto molto maggiore alla programmazione in grande.
	L'utilizzo di un linguaggio compilato porterebbe in questo senso ad un buon incremento delle prestazioni.
	 